\chapter{Introduction}
\label{sec:intro}

A pure functional language promises that a value, once made, never changes. Compiled to a
machine with mutable memory, the promise has a price: an update of a large structure must either
copy it or use a persistent representation that shares the unchanged parts. The price is
avoidable when nobody can tell the difference---when no part of the program will ever look at the
old version again. The long line of work on linear and uniqueness types
\citep{wadler1990linear,wadler1991use,smetsers1994guaranteeing,barendsen1996uniqueness,devries2008simplified},
on destructive-update and sharing analysis
\citep{bloss1989update,jones1989compile,wand2001set,cann1990sisal,aspinall2008usage,mazur2004thesis},
and more recently on reference counting with reuse
\citep{ullrich2019beans,reinking2021perceus,lorenzen2023fp2} and on modes in OCaml
\citep{lorenzen2024oxidizing} all ask, in one form or another, when an in-place write is
unobservable.

This thesis asks the question for one language and one data type, with an unusually strict
answer in mind. beni \citep{beni2026} is a strict, pure, Elm-like language that compiles to
JavaScript and is aimed at the browser. Its single sequence type, \code{List}, is backed by
JavaScript arrays. We want every \code{List} to be a \emph{plain} JavaScript array, written
\emph{in place} by every in-place edit the program asks for, and we want the compiler to guarantee
that this is safe. Concretely:

\begin{quote}
\emph{Editing a list in place while anything else still holds the old version is a compile error.}
Ownership is inferred, user code carries no annotations, and an accepted in-place edit writes its
array, always---in development and release builds alike---with no reference count, no run-time test
and no silent fallback to a copy.
\end{quote}

An in-place edit is a call of one of the core library's named list writers, such as
\code{List.set} or \code{List.push}. Building a list---a literal, a spread, a concatenation with
\code{++}, \code{List.map}---allocates a new array, exactly as the same expression written in
JavaScript would; the rule is about the writers, and \S\ref{sec:model-demands} says why syntax was
left out of it.

Two things make this goal different from most prior work. First, the result is an \emph{error},
not an optimisation. Run-time schemes such as Perceus \citep{reinking2021perceus}, Lean~4's
borrow inference \citep{ullrich2019beans} and Roc's in-place mutation of uniquely referenced
values \citep{teeuwissen2023roc}
update in place when a reference count says the value is unique and copy otherwise; the program is
always accepted, and its cost depends on facts the programmer cannot see. Static schemes that
\emph{optimise}---Sisal's copy elimination \citep{cann1990sisal}, Wand and Clinger's set
constraints \citep{wand2001set}, Mazur's compile-time garbage collection for Mercury
\citep{mazur2004thesis}---also fall back to a copy or to an unoptimised version. We want the opposite
contract: a program whose in-place edit would need a hidden copy is rejected, with a message that
names the holder of the old version and offers the explicit copy. Second, the language has no
references, no lifetimes and no ownership vocabulary in its surface syntax, so everything must be
inferred; the type-based uniqueness systems of Clean \citep{barendsen1996uniqueness,plasmeijer2011clean}
and Futhark \citep{henriksen2017futhark} express uniqueness as attributes in types; in Futhark the
programmer writes them on every function that consumes an argument.

\paragraph{Why it matters in a JavaScript target.}
A persistent sequence in JavaScript is not free. beni's current representation (\S\ref{sec:bg-lists})
is an array that switches to a 32-way radix trie when shared writes grow large; it costs runtime
code that every program that edits a list must ship, a representation test on every read site that
can meet a trie, and a conversion whenever a value crosses into JavaScript. In a prototype study
(\S\ref{sec:bg-measure}; method in Appendix~\ref{app:method}) in which we applied by hand the
transformations a static uniqueness pass would make to compiled beni code, building an array in a
fold by in-place writes ran at 0.14--0.47 times the time of the persistent representation, the
steady-state ticks of a million-cell grid at 0.44--0.48 times, and the runtime support for lists
shrank from 1\,514 to 706 bytes after compression once the trie was gone. The same study found the
obstacle: a UI runtime that keeps the lists it rendered makes every list in the model shared, so no
static rule could write them in place until the runtime gives up its hold on the old model when a
new one is computed---what we call \emph{handing the model over}.

\paragraph{Why beni can infer it.}
Three properties of the language do most of the work.
(1)~\emph{Strict, left-to-right evaluation in source order.} The order in which a read and a write
happen is fixed by the program text, so ``before the update'' is decidable from the text, as in
Futhark \citep{henriksen2022uniqueness} and in Smetsers et al.'s classification of argument
positions into ``preliminaries'' and groups of ``alternate arguments''
\citep[\S4]{smetsers1994guaranteeing}. Care is needed with what ``before'' means: an argument
written before an update may still be \emph{held}, unread, until a later call reads it
(\S\ref{sec:model-unique}).
(2)~\emph{No automatic currying.} Every call in beni is saturated; partial application is written
as an explicit lambda. A partial application that \emph{silently} captures a unique argument---a
case Clean, de Vries et al.\ and OxCaml all handle with special rules
\citep{barendsen1996uniqueness,devries2008simplified,lorenzen2024oxidizing}---cannot be written.
OxCaml's designers note that the complications of currying for modes ``go away if we imagine a
change to the language forbidding currying'', since the explicit closure ``would infer its own result
mode'' \citep[\S6.4]{lorenzen2024oxidizing}. What remains is the explicit closure, and it is not
free: a closure that is returned from a function, stored in a structure or captured by another
closure carries its captured arrays, and the fact that it may be called only once, to places the
lambda's text does not show. Function types therefore carry inferred capture and affinity facts,
close to OxCaml's modes on arrows, with no annotation (\S\ref{sec:model-functions}).
(3)~\emph{An existing inference of per-function facts.} beni's type checker already infers two
effect bits for every function type, solves them as a directed constraint graph after each module,
and publishes them, hashed, in the module interface. Ownership facts can ride in the same places: a
few more facts per function-type site, the same inclusion edges, the same interface block, the same
rule for when importers must be re-checked.

\section{The thesis on one page}
\label{sec:intro-onepage}

\paragraph{The property.} An in-place edit is accepted when its operand's array is \emph{unique at
the site}: every other holder of that array is dead at that point. A holder is dead when no later
step reads through it. A dereference that \emph{completes} before the edit is therefore never
sharing, however many there are; a value merely held across the edit---a variable waiting to be
passed to a later call, a field of a record still being built, a closure called afterwards---is a
live holder. This is Boyland's alias burying, under which aliases may exist provided they ``all be
dead'' \citep[\S3.2]{boyland2001alias}, and Rust's liveness-based borrows \citep{rfc2094}, applied to
a language with no references. It is exactly what a plain array needs: the write is invisible if and
only if nobody reads the old version afterwards.

\paragraph{The analysis.} Per function, one forward pass over the body in evaluation order computes,
for each variable and path, the set of arrays it may hold; paths reach into records, tuples,
constructors, list elements and the environments of closures. The flow is directional, in the
style of Andersen's points-to analysis, never unification: Emre et al.\ measured that four
necessarily unsafe pointers can poison 85\,\% of the raw pointers under a unification-based model
\citep[\S5]{emre2023aliasing}. A backward liveness pass per (variable, path) follows, and then, at
each update, the check. Calls use the callee's summary. A recursive group starts from the most
optimistic summaries and weakens to a fixpoint over a finite lattice, the search that Aspinall,
Hofmann and Kone\v{c}n\'y report finds the best typing for their first-order system
\citep[\S6.2]{aspinall2008usage} and the direction that Lean's borrow inference
\citep[\S5.2]{ullrich2019beans} and Brandon et al.\ \citep[\S4.4]{brandon2026borrows} take.

\paragraph{What it needs from the rest of the system.} The core library and the platforms
\emph{assert} summaries for their primitives, as their foreign declarations already assert effects;
those assertions, together with a handful of runtime contracts, are the trusted surface, and it is
larger than the list writers alone (\S\ref{sec:soundness-trusted}). Every function a UI runtime
calls on the program's behalf---\code{init}, \code{update}, \code{view}, subscriptions, the bodies of
commands, derived values, event listeners---has an entry protocol saying which arguments it receives
owned or only to read, and what the runtime keeps of its result (\S\ref{sec:hard-tea}). beni's
current browser runtime keeps every rendered list, so on that platform every in-place edit of a
rendered list is an error; an experimental platform that compiles the view away can host the rule
under the obligations stated there.

\paragraph{What it costs.} Every program that keeps an old version and then edits it in place---an
undo history, a ``before and after'' comparison, a test that holds its input, a command that captured
the list---is rejected until it writes \code{List.copy}. Persistent sharing becomes an explicit $O(n)$
copy, and the trie goes. A closure that consumes what it captured may be called at most once, so
\code{List.map xs (λx → List.push acc x)} is an error and a fold is the way. A body edit that changes
a public function's summary re-checks its importers. Because list syntax builds new arrays, an
accumulator written \code{acc ++ [ x ]} copies at every step, as it would in JavaScript; the checker
warns where a named writer would have been safe. And the two real applications we examined edit
their lists with \code{List.map}, \code{List.filter} and \code{++}: under the rule they contain no
in-place edit at all, so they are accepted and gain nothing (\S\ref{sec:disc-apps}).

\paragraph{The proposal.} Build the checker as a \emph{report-mode pass first}, together with a
differential run that executes every accepted program under both semantics with a trap on any read of
an out-of-date array; run both over the core library, the compiler's test programs, the two
applications and a corpus of programs written outside the project; classify every rejection as real
sharing or a limit of the checker; and decide against criteria fixed before the run
(\S\ref{sec:evaluation}). We believe the soundness argument has the right shape; the rule's value
depends on what no desk analysis can settle: the rejection rate on code people actually write, and
how much code would use the in-place writers at all.

\section{Contributions}

This thesis is a design proposal under evaluation. Its contributions are:
\begin{itemize}
\item A precise statement of the property to check: \emph{uniqueness at the update site, decided by
liveness} (\S\ref{sec:model-unique}). It is \emph{incomparable} with Clean's occurrence-based
uniqueness: it accepts reads before a write and dead aliases, which Clean's marking rejects, in the
spirit of Boyland's alias burying \citep{boyland2001alias} and Rust's two-phase borrows
\citep{rfc2025}; Clean's unique spines with unique elements accept element edits that our abstraction
accepts only when it can prove a container's contents exclusive (\S\ref{sec:model-excl}).
\item An inference algorithm (Chapter~\ref{sec:inference}): a directional, flow-sensitive,
field-sensitive cell analysis whose paths reach into closure environments, a per-path liveness pass
over an evaluation-order normal form, four side conditions at each update (\Known, \Unescaped,
\Dead, \Disjoint), affinity for closures that consume what they capture, function summaries
polymorphic over the ownership behaviour of function-typed parameters, an optimistic fixpoint over a
stated finite lattice for recursive groups, and publication of summaries in module interfaces.
\item A soundness argument (Chapter~\ref{sec:soundness}): under twelve stated assumptions about the
runtime, the core library and the platforms, an accepted program produces the same observable
trace under value semantics and under in-place semantics. We give the theorem, the lemmas and proof
sketches, the trusted surface, the places we consider weakest, and twelve programs that an earlier
version of the rules wrongly accepted, each traced through the revised rules.
\item An entry protocol for every function a TEA (Elm Architecture) runtime calls, and a treatment of
the hard cases (Chapter~\ref{sec:hard}): closures, nested containers, generic code, fibers and
commands, the foreign-function boundary, views, mutable cells.
\item An error-message design (Chapter~\ref{sec:errors}) that names three program points, writes the
fix, and tells \emph{real sharing} from \emph{a limit of the checker} by a derivation tag carried
on every fact, together with a stability rule for the accepted set across compiler releases.
\item An account of the trade-off (Chapter~\ref{sec:discussion}), the language and library changes the
rule needs or invites (Chapter~\ref{sec:disc-options}), and an evaluation plan
(Chapter~\ref{sec:evaluation}) whose first experiment tests soundness and whose decision criteria are
stated in advance. \emph{No results of that evaluation exist yet}; we report none.
\end{itemize}

\section{How to read this thesis}

Chapter~\ref{sec:background} introduces beni, its runtime and its compiler at the depth the
argument needs, and states the question precisely. Chapter~\ref{sec:model} is the model,
Chapter~\ref{sec:inference} the inference, Chapter~\ref{sec:soundness} the soundness argument,
Chapter~\ref{sec:hard} the hard cases, Chapter~\ref{sec:errors} the error messages, and
Chapter~\ref{sec:compiler} the integration with the rest of the compiler.
Chapter~\ref{sec:discussion} states what the rule costs, Chapter~\ref{sec:disc-options} the
language changes it needs or invites, Chapter~\ref{sec:evaluation} the experiment to run before
deciding, Chapter~\ref{sec:related} related work, and Chapter~\ref{sec:conclusion} the conclusions,
recommendations and open questions. Appendix~\ref{app:examples} works programs through the rules;
Appendix~\ref{app:priorart} records, line of work by line of work, what the design takes from the
prior art and what it declines; Appendix~\ref{app:method} gives the method behind every number we
quote from our own measurements.

\section{Terms}
\label{sec:intro-terms}

The following terms are used throughout; others are defined where they first appear.

\begin{description}[leftmargin=1.2em,font=\normalfont\bfseries]
\item[Cell.] One JavaScript array that backs a \code{List} value. Two \code{List} values that are
the same cell are \code{===}; editing a cell in place changes what every reference to it sees. A
closure that may be called only once is treated as a cell too (\S\ref{sec:model-functions}).
\item[Holder.] Anything from which a cell can be reached and read later: a variable binding, a
field of a record, an element of another list, a closure's environment, an argument evaluated but
not yet passed, a fiber's saved locals, a command's thunk, the runtime's own state, a JavaScript
file, a mutable cell.
\item[Content holder, identity holder.] A content holder may read the cell's elements; an identity
holder only ever asks \code{a === b} of it. Only content holders can observe an in-place edit;
identity holders are a separate hazard (\S\ref{sec:bg-question}).
\item[Dead.] A holder is dead at a point when no execution continuing from that point reads through
it. A dead holder may exist; it just never looks.
\item[Unique at a site.] A cell is unique at an update site when every holder of it is dead once the
update's arguments have been evaluated---the operand's own binding included. This is the property
the checker proves (\S\ref{sec:model-unique}).
\item[Demand.] A function whose summary says it writes a parameter in place: the core library's named
writers, such as \code{List.set} and \code{List.push}. Nothing in user code is a demand; user code
only inherits demands by calling them.
\item[Summary.] What the checker infers per function and publishes in its interface: for each
parameter path, whether the function ignores it, reads it, lets its result alias it, lets something
unseen keep it, or writes it; for each result path, which parameters it may alias; for each
function-typed parameter, how often it is called and what is passed to it
(\S\ref{sec:model-summary}).
\item[Hand-over.] A UI runtime hands the model over when, after calling \code{update}, it keeps no
reference through which it could read a list of the old model (\S\ref{sec:hard-tea}).
\item[Direct platform, derived value, write set.] An experimental browser back end of beni that
compiles the view away; a derived value is what a binding inside \code{view} becomes there; and the
write set of a message is the set of model paths its \code{update} branch may change, computed by a
separate whole-program analysis (\S\ref{sec:bg-browser}).
\item[Same-shape rebuild.] A library function such as \code{map}, \code{filter} or \code{sort} whose
result could reuse its input's array; whether in-place forms of these belong in the library is
Change~1 of Chapter~\ref{sec:disc-options}.
\item[False error.] A rejection of a program that would in fact have run correctly in place. The
checker is designed to be sound and incomplete; if the argument of Chapter~\ref{sec:soundness} holds,
every imprecision is a false error, never a wrong answer. Chapter~\ref{sec:errors} says how the
message tells the two apart.
\item[Value semantics, update semantics.] The meaning beni has today (every update builds a new list;
the old one is untouched) and the meaning the checker licenses (an accepted update writes the cell).
The soundness claim is that the two cannot be told apart (Chapter~\ref{sec:soundness}).
\end{description}
